\documentclass[11pt,letterpaper]{article}
\usepackage[margin=0.7in,headheight=15pt]{geometry}
\usepackage[T1]{fontenc}\usepackage{lmodern,amsmath,amssymb,fancyhdr,parskip,hyperref}
\hypersetup{colorlinks=true,urlcolor=blue}
\pdfinfoomitdate=1\pdftrailerid{}\pdfsuppressptexinfo=15
\pagestyle{fancy}\fancyhf{}\lhead{Week 84 / Public knowledge}\rhead{Facilitator draft}\cfoot{Bellingham Math Circle / Week 84 / facilitator / \thepage}
\setlength{\parskip}{7pt}\setlength{\parindent}{0pt}
\begin{document}
\section*{Mathematical overview}
Knowledge means that every remaining possibility compatible with a player's private view gives the same answer. In the three-person cafe, each knows only their own YES/NO preference and earlier public responses to ``Does everyone want a drink?'' A personally NO player can answer NO immediately. A first ``not enough information'' therefore proves that first preference is YES. Later NO need not mean the later speaker personally says NO.

For three hidden-to-self blue/red headbands, announce publicly that at least one is blue. With ideal evidence-complete reasoners reporting simultaneously, if there are $k$ blue headbands, precisely those blue dolls can prove blue on round $k$, for $k=1,2,3$. Induction uses a sole blue's immediate knowledge, then the public absence of earlier blue declarations. A red doll who knows red still reports ``not yet'' because only blue-proof is requested.

Experiments use complete movable possibility cards. The all-worlds check, not a child's confidence or delay, establishes a declaration. Grades 2--3 enter through cafe; Grades 3--5 coach dolls. Optional Grades 4--5 finite number grid uses a \emph{bounded} domain with both boundaries. No transfinite dialogue appears on student pages.
\newpage
\section*{Setup, information and common launch}
Use three named dolls and B/R headbands, cafe YES/NO cards, public response cards and eight three-position configuration cards. Print cafe and headband configuration cards (pages 4--5), reply cards (pages 6--7), and page 8's cuttable staging labels and three response slots. Arrange all eight 3.45 by 1.85 inch configuration cards openly in a four-by-two tabletop display; allow about 14.55 by 3.95 inches including gaps. The labels mark open tabletop areas, not card containers. Name/order labels must match throughout. A doll sees the other two colors but cannot see its own. Children coach dolls and may inspect a full configuration only as facilitators/spectators; their coaching must use the doll's permitted information. The actual arrangement stays fixed until a fresh game.

For the cafe, the player privately knows their own preference and can hear previous answers; the audience's possibilities are shared separately. For hats, retain all eight arrangements before the public announcement; remove only RRR at the announcement. Do not secretly tell the exact blue count.

Use a large tabletop display for the seven or eight configuration cards; the printed strips label its areas rather than containing the full set. All cards must remain visible. Launch after materials play: temporarily point to or group the possibilities compatible with one doll's private observation, using page 2's two-doll BR worked example. Put them back afterward. In an actual BBR deal, doll 1 considers BBR/RBR, doll 2 considers BBR/BRR, and doll 3 considers BBR/BBB; these are three different private views of the same seven-world public set. Private matching is temporary; it does not remove configurations from the shared public set. Compare the doll's information with the audience's. Set response cards face down; reveal every doll's choice together. Keep one unchanged public set throughout each round, then remove incompatible configurations together. Sequential elimination inside a round leaks extra information.

One anchored adult is needed per table; for the upper triplet rotate two coaches and the recorder/referee, keeping all three doll identities. First route: cafe 10 minutes; one-blue hats 10; two-blue 15; compare three-blue if ready, then share. Optional grid is an encore. Reading/recording help is routine; if adults choose every elimination, the mathematical agency has shifted and should be recorded.
\newpage
\section*{Cafe solutions and headband proof}
Problem 1's four conversations leave, in order: YYN; NYY/NYN/NNY/NNN; YNY/YNN; YYY. The second set retains any preferences for the final two speakers, since the first N has already established the answer NO. In the third conversation first uncertainty fixes first Y, then second NO fixes second N; the third preference remains unknown. Cafe example Y,Y,N: first uncertainty removes every arrangement with first N. Second uncertainty removes every arrangement with second N among the survivors. The third says NO, so only YYN remains. By contrast NYY starts with NO; the later preferences remain unidentified. If the first says NO, a later YES-preferring speaker already knows the answer NO and also says NO. Do not infer every NO speaker has preference N.

Problem 2, hats with one blue: that doll sees two reds; since RRR was publicly excluded, its own color must be blue, so it declares on round 1. A red doll sees a blue and still has both own-color possibilities initially.

Hats with two blues: each blue sees one blue, one red. If it were red, the visible blue would be sole blue and declare at round 1. Nobody declares in the actual configuration, so both blue dolls infer blue and declare on round 2. The red doll sees two blues; after round 1 its red/two-blue and blue/three-blue alternatives remain.

Hats with three blues: each sees two blues. If it were red, those two would declare at round 2. Silence through round 2 removes that alternative; all three declare round 3. A blue proof must be checked against possibilities \emph{before} simultaneous revelation, not the information created by a neighbor revealing earlier.

Problem 3: on round 1, the three one-blue configurations have distinct response patterns (the sole blue proves). BBR, BRB, RBB and BBB all give three not-yets and stay grouped. On round 2 the three two-blue configurations give their three distinct two-proof patterns, while BBB still gives three not-yets. Thus all seven deals have distinct public response histories by round 2, although BBB gives its three blue proofs only on round 3.

Problem 4: without the announcement, RRR remains possible. In every private view each doll has an own-red alternative, so nobody can prove blue in round 1 in any configuration. Public all-not-yets removes nothing, and induction repeats this unchanged state in every round. No blue proof ever follows from this protocol alone.

The finite inference rule treats all agents as perfect and truthful. Real hesitation, mistakes and different processing times are not data. If children disagree, return to visible possibility cards instead of judging their reasoning speed.
\newpage
\section*{Optional bounded hidden-number grid}
Numbers are different members of $\{0,1,\ldots,6\}$. Alice knows her row; Bob knows his column; everybody knows this domain. Remove the diagonal first. For a worked deal Alice=3, Bob=2, alternate reports on who has the larger number. This is a finite adaptation, not a truncation retaining the answer of Hamkins's unbounded puzzle.

Alice's first uncertainty removes rows 0 and 6: either endpoint knows who is larger. Bob's uncertainty removes columns 0,1,5,6, given remaining Alice rows 1--5. Alice's next uncertainty leaves only row 3, since rows 1,2 know Bob is larger and rows 4,5 know Alice is larger. The public possibilities are $(3,2),(3,4)$; Bob with private 2 now knows both numbers. The audience does not yet know Bob's number. Announcing ``Alice is larger'' finally leaves $(3,2)$.

If the finite grid is absent from the selected student draft, conduct this only as a later adult-led whiteboard encore. A child need not copy or complete 49 cases individually; share a board or use counters. Distinguish knowing the comparison, knowing both numbers, and publicly naming them.

Changing the domain to 0--5 changes the elimination pattern because of the upper boundary. It is a worthwhile investigation once the original protocol is controlled. Statements about all possible finite conversations in transfinite puzzles are additional information, not merely extra waiting.
\newpage
\section*{Provenance, return visits and observation record}
Sources: Joel David Hamkins's epistemic-logic presentation (Oxford October 2019) contains the classic Logic Cafe and blue-eyed islanders mechanisms. \href{https://jdh.hamkins.org/i-know-that-you-know-that-i-know-that-you-know-oxford-october-2019/}{Oxford epistemic-logic presentation (October 2019)}. Classic puzzles are presented there, not claimed here as his inventions. Our three-doll version, possibility cards and finite bounded grid are new adaptations.

Teacher inspiration only: \url{https://jdh.hamkins.org/now-i-know/} (2015), and ``Cheryl's Rational Gifts'' with its separate solution. These use stronger announcements about hypothetical continuations and genuinely transfinite structure. They do not justify asking elementary children to simulate infinitely many rounds.

Encore: children invent two hidden cafe arrangements giving the same public conversation, or compare simultaneous hat reports with a deliberately changed sequential protocol. Always state changed information rules before analyzing a new game. Existing Week 6 code breaking asks about information but does not already implement higher-order public knowledge.

Pedagogy: Rozhkovskaya, Lesson 3 ``At the lesson,'' item 1 (part0013) describes children trying a logic puzzle before the instructor makes possibilities visible in a table. We replace much of the table writing by supplied movable configuration cards. Lesson 8 item 2 supports lighter transcription; these proposals remain untested.

Unpiloted. US Letter, single-sided, Actual Size. Doll/headband fit, opaque backs, card cutting, timing and classroom procedure have not been rehearsed. Record actual eliminations children controlled, rounds reached, accidental information leaks, misunderstandings of private versus public knowledge, and questions for another visit.
\end{document}
